\section{Introduction}
\label{sec:intro}

Knowing what the sky is doing above a point still takes a trained observer or a ceilometer. An observer reports
cloud genus, amount and base height by eye following the WMO conventions~\cite{wmo2017atlas}; a ceilometer
measures the base height along one vertical line and costs more than most of the stations that would like one.
Upward-looking cameras cost a few thousand rupees, and a decade of work has shown that cloud type, cover and even
base height can be read from their pictures~\cite{zhang2018cloudnet,dev2017colorbased,liu2020mgcd}. Yet almost
none of that work transfers to a camera it was not trained on, and much of it reports numbers that do not survive
a second look.

Three problems stand between the published results and a deployable model. First, \emph{leakage}: the popular
recognition sets are built from web images and time series; exact and near-duplicate images straddle train and
test splits, and the same scene appears minutes apart. Our audit (\cref{sec:data}) finds that the 91--97\,\%
accuracies reported on CCSN~\cite{zhang2018cloudnet} fall to 57--62\,\% once duplicates are grouped. Second,
\emph{forced labels}: experts disagree on cloud type on a large fraction of images; a single hard label hides a
measurable ceiling, and a model scored against it is scored against noise. Third, \emph{one camera}: a classifier
that has never seen another lens, exposure or horizon does not know that it is looking at a new sky; it answers
anyway.

This paper treats those three problems as the method. We assemble seven public datasets and our own camera into a
single corpus with grouped, day-blocked and leave-one-dataset-out splits, WMO-harmonised \emph{set-valued} labels
and a human ceiling measured from five annotators; we pre-register the success criteria and lock the test sets
before any model is trained, and read them once. On that footing we build STRATIA (\cref{fig:teaser}): a frozen
DINOv3 ViT-S/16 backbone~\cite{simeoni2025dinov3}, a \emph{SkyRay} encoder that adds to every token the direction it
looks at (from the camera calibration) and where the Sun is, with geometry dropout so that ``unknown'' geometry is a
valid input, and heads for genus and \'etage sets, ordinal cover, a sky/cloud/glare parse with cloud layers, an
ordinal cloud-base height distribution supervised by ceilometers of an all-sky camera network~\cite{eye2sky}, and
an out-of-distribution score that turns into an explicit \emph{uncertain} state. The whole model exports to one ONNX
graph with its calibration folded in, and runs inside an open desktop application on a laptop GPU.

\paragraph{Contributions.}
\begin{itemize}
\item A leakage-audited, WMO-harmonised benchmark protocol for sky cameras: grouped and temporal splits, LODO
  folds, a held-out station, set-valued labels, a measured expert ceiling, and locked test sets with a logged
  single read (\cref{sec:data}).
\item STRATIA: geometry-aware heads on a frozen foundation backbone, trained on cached features in stages, with
  temperature and split-conformal calibration and an OOD score (\cref{sec:method}).
\item Locked-test results with bootstrap intervals (\cref{sec:experiments}): expert-level agreement on
  \'etage and cover, ceilometer-level cloud-base height that transfers to a site never seen in training, and two
  pre-registered negatives reported as such.
\item A few-shot adaptation to the authors' own camera from three afternoons, verified on a locked fourth day, and
  a working deployment (\cref{sec:b0268}).
\end{itemize}
